\subsection{Total elimination depth}
\label{subsec:np-total-depth}

Auxiliary pairs encode the normalized parity source of
Theorem~\ref{thm:xor-source} in a tripartite graph. Completing its parts
to cliques follows Bonnet's construction~\cite[Section~3.1]{BonnetTreewidth}.
Three added independent sets then transfer the number of satisfied
equations to total depth, giving the limiting gap ${189/186=63/62}$.

\subsubsection{The three-path objective}

Let $G$ be a simple tripartite graph with specified nonempty independent
parts $A\mathbin{\dot\cup}B\mathbin{\dot\cup}C$, and put $N=|V(G)|$.
Complete each part to a clique by setting
\[
 \widehat G=\left(V(G),\ E(G)\cup
                 \bigcup_{X\in\{A,B,C\}}\binom{X}{2}\right).
\]
For an integer $M\geq1$, form $H^\Sigma_M(G)$ by adding disjoint
independent sets $T_A,T_B,T_C$, where $|T_X|=M$ and
$N_{H^\Sigma_M(G)}(t)=X$ for every $X\in\{A,B,C\}$ and $t\in T_X$.
The resulting graph is simple and unweighted.

In an elimination forest $F_G$ of $\widehat G$, the vertices of each
clique part lie on one ancestor path. Write $y_X$ for the deepest vertex
of part $X$, with depths starting at one, and define
\[
 W(F_G)=\sum_{X\in\{A,B,C\}}
          \bigl(\operatorname{depth}_{F_G}(y_X)+1\bigr),
 \qquad W(G)=\min_{F_G}W(F_G).
\]
Here and below, minima over $F_G$ range over all elimination forests
of $\widehat G$, and $\totd(F_G)$ denotes the sum of their vertex depths.
The following normalization and path identities are proved in
Appendix~\ref{app:total-depth-tools}.

\begin{lemma}[Normalization of added vertices]
\label{lem:np-total-depth-added-vertices}
Every elimination forest of $H^\Sigma_M(G)$ can be transformed without
increasing total depth so that all vertices of each $T_X$ are leaf
children of the deepest vertex of $X$. Consequently,
\begin{equation}\label{eq:np-total-depth-added-vertices-identity}
 \totd(H^\Sigma_M(G))
 =\min_{F_G}\bigl(\totd(F_G)+MW(F_G)\bigr),
\end{equation}
and
\begin{equation}\label{eq:np-total-depth-added-vertices-bounds}
 MW(G)
 \leq\totd(H^\Sigma_M(G))
 \leq MW(G)+\frac{N(N+1)}2.
\end{equation}
\end{lemma}

For $X\in\{A,B,C\}$, let $Y,Z$ be the other two parts. Let
$S_0\subseteq V(G)$ range over covers of all $X$--$(Y\cup Z)$ edges, and let
$S_1\subseteq(Y\cup Z)\setminus S_0$ be a minimum vertex cover of
the remaining $YZ$ graph. Define
\begin{equation}\label{eq:np-total-depth-cover-cost}
 \operatorname{cost}_X(G)=\min_{S_0}
 \left(2|S_0|+\vc\bigl(G[(Y\cup Z)\setminus S_0]\bigr)\right).
\end{equation}
All cover graphs refer to $G$, before clique completion. In the
associated forest, $S_0$ lies on a path common to all three parts,
$S_1$ lies on a path shared by $Y,Z$, and the remaining vertices of
each part lie on its own path.

\begin{lemma}[Cost of the three paths]
\label{lem:np-total-depth-path-cost}
For every tripartite graph $G$ with nonempty specified parts,
\begin{equation}\label{eq:np-total-depth-path-cost-formula}
 W(G)=N+3+\min\{\operatorname{cost}_A(G),
                 \operatorname{cost}_B(G),\operatorname{cost}_C(G)\}.
\end{equation}
\end{lemma}

\subsubsection{The tripartite parity construction}

Fix $0<\epsilon<1/4$ and let $F$ be an instance from
Theorem~\ref{thm:xor-source}, with variable set $\mathcal X(F)$,
$m\geq1$ equations, and satisfying
\eqref{eq:xor-normalization}. Thus $d_x$ is the number of equations
containing variable $x$, $d_x\in6\mathbb Z_{>0}$, and
$\sum_{x\in\mathcal X(F)}d_x=3m$. The source occurrence bound is $d_x\leq b_\epsilon$.
Construct $G=G_F^\Sigma$ with independent parts $A,B,C$ as follows.
All blocks and all auxiliary vertices mentioned below are pairwise
disjoint. An \emph{auxiliary pair} consists of fresh vertices
$b\in B$, $c\in C$ and the edge $bc$. Attaching an endpoint to a
block means adding every edge between that endpoint and the block.

\begin{description}
\item[Literal blocks.]
 For every $x\in\mathcal X(F)$, put two blocks $A_x^0,A_x^1$ in $A$, with
 $|A_x^0|=|A_x^1|=9d_x/2\in\mathbb Z_{>0}$.
\item[Opposite pairs.]
 For every $x\in\mathcal X(F)$, create $2d_x$ auxiliary pairs. For $d_x$ of them,
 attach $b$ to $A_x^0$ and $c$ to $A_x^1$; for the other $d_x$,
 attach $b$ to $A_x^1$ and $c$ to $A_x^0$.
\item[Equation pairs.]
 In each equation $e\in\{1,\ldots,m\}$, order its distinct variables as $x_1,x_2,x_3$
 and write $a_e$ for its right-hand side. For every satisfying local
 assignment $\sigma=(\sigma_1,\sigma_2,\sigma_3)\in\{0,1\}^3$
 with $\sigma_1+\sigma_2+\sigma_3=a_e\pmod2$, and every
 $j\in\{1,2,3\}$, create one auxiliary pair. Attach $b$ to all
 three blocks $A_{x_i}^{1-\sigma_i}$, $i\in\{1,2,3\}$, and
 attach $c$ to $A_{x_j}^{\sigma_j}$. This gives twelve pairs per
 equation.
\item[Padding pairs.]
 Add $18m$ auxiliary pairs having no neighbors in $A$.
\end{description}
Add no other edges to $G$. In particular, its $BC$ graph is a matching.
There are $6m$ opposite pairs, $12m$ equation pairs, and $18m$ padding
pairs, so
\begin{equation}\label{eq:np-total-depth-part-sizes}
 |A|=27m,\qquad |B|=|C|=36m,\qquad N=99m.
\end{equation}
The two vertex types satisfy
\[
 \deg_G(v)=2d_x+8d_x=10d_x
 \quad(v\in A_x^0\cup A_x^1),\qquad
 \deg_G(v)\leq3(9b_\epsilon/2)+1
 \quad(v\in B\cup C).
\]
The second bound accounts for at most three literal-block
neighborhoods and the paired endpoint. Since $b_\epsilon\geq6$,
\begin{equation}\label{eq:np-total-depth-degree}
 \Delta(G_F^\Sigma)\leq14b_\epsilon.
\end{equation}
This is the degree before completing the parts and adding $T_A,T_B,T_C$.

Write $F_0$ for the initial parity source, $G=G_F^\Sigma$, and
$H=H^\Sigma_M(G)$. The reduction has four stages:
\[
 F_0\xrightarrow{\text{normalize}}F
 \xrightarrow{\text{encode}}G
 \xrightarrow{\text{complete}}\widehat G
 \xrightarrow{\text{add }T_A,T_B,T_C}H.
\]
The source is from~\cite[Theorem~10]{HazanSafraSchwartz2003}, with
normalization in Appendix~\ref{app:xor-normalization}.
Proposition~\ref{prop:np-total-depth-parity-identity} analyzes the
encoding, the clique completion follows~\cite[Section~3.1]{BonnetTreewidth},
and Lemma~\ref{lem:np-total-depth-added-vertices} handles the added sets.

\begin{proposition}[The encoded parity optimum]
\label{prop:np-total-depth-parity-identity}
The tripartite graph above satisfies
\begin{equation}\label{eq:np-total-depth-parity-W}
 W(G_F^\Sigma)=192m-6\OPT(F)+3.
\end{equation}
In particular, $W(G_F^\Sigma)\geq186m+3$.
\end{proposition}

\begin{proof}
We compute the $A$-first cost and lower-bound the two competing costs
in~\eqref{eq:np-total-depth-path-cost-formula}.

\emph{The $A$-first cover.}
By Lemma~\ref{lem:total-depth-cover-normalization}, each literal block
can be taken wholly in $S_0$ or wholly outside it. Omitting a literal
block forces all its auxiliary neighbors into $S_0$. Retain exactly
these forced endpoints: removing any unforced auxiliary endpoint from
$S_0$ saves two and increases the residual cover cost by at most one.
The residual $BC$ graph remains a matching, so each pair contributes
one when neither endpoint is forced, two when one is forced, and four
when both are forced. Relative to an unforced pair, the extra costs
are therefore $0,1,3$.

With every literal block in $S_0$, the cost is
\begin{equation}\label{eq:np-total-depth-A-baseline}
 2|A|+36m\overset{\eqref{eq:np-total-depth-part-sizes}}{=}90m.
\end{equation}
For variable $x$, omitting one block saves $9d_x$
and incurs opposite-pair extra cost $2d_x$, giving net saving $7d_x$.
Omitting both saves $18d_x$ and incurs extra cost $6d_x$, giving net
saving $12d_x$. Allocate these net savings equally among the $d_x$
equations containing $x$.

For an equation $e$, let $s$ be the number of its variables with
exactly one omitted block and let $b$ be the number with both blocks
omitted. A variable with exactly one omitted block $A_x^v$, where
$v\in\{0,1\}$, is fixed to value $v$. Call a satisfying local assignment
compatible if it agrees with every omitted literal, and let $n_e$ count such assignments.
Thus $n_e=0$ when $b>0$, since the two omitted literals of one
variable prescribe conflicting values.
Write $\Delta_e$ for the allocated variable savings minus the extra
costs of the twelve equation pairs. Then
\begin{equation}\label{eq:np-total-depth-local-saving}
 \Delta_e=3s+4b-12+(3+s)n_e.
\end{equation}
To count the extra equation cost, note that each omitted literal
forces an endpoint in $C$ in exactly two of the four satisfying assignments.
There are therefore $2s+4b$ forced endpoints in $C$ in total.
A compatible assignment forces no endpoint in $B$ and exactly $s$
in $C$. An incompatible assignment forces all three endpoints in
$B$, at extra cost three, and each forced endpoint in $C$ then
costs two more. Hence the extra equation cost is
\[
 E_e=n_es+3(4-n_e)+2(2s+4b-n_es)
 =12+4s+8b-(3+s)n_e.
\]
Subtracting this from the allocated savings $7s+12b$ proves
\eqref{eq:np-total-depth-local-saving}.
For $b=0$ and $s\leq2$, the unfixed variables give
$n_e=2^{2-s}$. Substitution in~\eqref{eq:np-total-depth-local-saving}
and the bound $s+b\leq3$ therefore give
\begin{equation}\label{eq:np-total-depth-saving-table}
\begin{gathered}
 \Delta_e=-s(3-s)/2\leq0 \quad (b=0,\ s\leq2),\qquad \Delta_e=3s+4b-12
  \overset{s+b\leq3}{\leq}b-3\leq0 \quad (b>0),\\
 \Delta_e=6n_e-3\in\{-3,3\} \quad (b=0,\ s=3).
\end{gathered}
\end{equation}
In particular, an equation with an unfixed variable has saving at most
zero. Complete the fixed values by independent fair bits for variables
having neither or both blocks omitted. By
Lemma~\ref{lem:partial-parity-completion}, each unfinished equation
is satisfied with probability $1/2$, so its expected saving after
completion is $(3-3)/2=0$. Fully fixed equations keep their saving.
Writing $\Delta_e^{\rm old}$ and
$\Delta_e^{\rm new}$ for the respective savings, linearity gives
\[
 \mathbb E\!\left[\sum_{e=1}^m\Delta_e^{\rm new}\right]
 \overset{\text{linearity}}{=}
 \sum_{e=1}^m\mathbb E[\Delta_e^{\rm new}]
 \overset{\eqref{eq:np-total-depth-saving-table}}{\geq}
 \sum_{e=1}^m\Delta_e^{\rm old}.
\]
No independence between equations is needed. Some outcome attains at
least its expectation. A full assignment satisfying
$k$ equations has saving $3k-3(m-k)=6k-3m$. Conversely, any assignment
is realized by omitting the block matching each assigned value. If
$S_{\max}$ denotes the maximum total saving over the normalized block
choices, this proves
\begin{equation}\label{eq:np-total-depth-max-saving}
 S_{\max}\overset{\eqref{eq:np-total-depth-saving-table}}{=}6\OPT(F)-3m.
\end{equation}
Therefore
\begin{equation}\label{eq:np-total-depth-A-cost}
 \operatorname{cost}_A(G)
 \overset{\text{Lem.~\ref{lem:total-depth-cover-normalization}},\,\eqref{eq:np-total-depth-A-baseline}}{=}
   90m-S_{\max}
 \overset{\eqref{eq:np-total-depth-max-saving}}{=}
   90m-(6\OPT(F)-3m)=93m-6\OPT(F).
\end{equation}

\emph{The competing covers.}
Consider a $B$-first or $C$-first cover. Normalize each literal block
using Lemma~\ref{lem:total-depth-cover-normalization}. Its state is
$T$ if it is in $S_0$, $U$ if it is in $S_1$, and $V$ if it is
in neither. For a block of variable $x$, these states cost
$9d_x,9d_x/2,0$, respectively.

An \emph{ordinary pair} has its $b$ endpoint attached to one literal
block, called its tail, and its $c$ endpoint attached to one literal
block, called its head. In the following indicators, $\text{tail}$
and $\text{head}$ denote the states of these blocks. For fixed block
states, its minimum auxiliary cost in a $B$-first cover is
\begin{equation}\label{eq:np-total-depth-rival-pair}
 2+\mathbf1[\text{tail}\ne T,\ \text{head}=V].
\end{equation}
The middle edge forces at least one endpoint into $S_0$, at cost two.
If the tail is $T$, choosing $c\in S_0$ attains this cost. Otherwise
$b\in S_0$ is forced. The head attachment is then already covered
when the head is $T$ or $U$; when it is $V$, the cheapest choice is
to add $c$ to $S_1$, costing one more. In a $C$-first cover the same
formula holds with tail and head exchanged. Here $\mathbf1[\mathcal E]$
is one if condition $\mathcal E$ holds and zero otherwise.

For an equation pair, index the three blocks adjacent to its $b$
endpoint by $i\in\{1,2,3\}$, write $\text{tail}_i$ for their states,
and write $\text{head}$ for the state of the block adjacent to $c$.
In the $B$-first case its extra cost above two is
$\mathbf1[\text{head}=V]\max_{1\leq i\leq3}\mathbf1[\text{tail}_i\ne T]$.
In the $C$-first case it is
$\mathbf1[{\text{head}\ne T}]\max_{1\leq i\leq3}\mathbf1[\text{tail}_i=V]$.
Each maximum dominates the average of its three entries; in particular,
\[
 2+\mathbf1[\text{head}=V]\max_{1\leq i\leq3}
                        \mathbf1[\text{tail}_i\ne T]
 \overset{\max\geq\text{average}}{\geq}
 \frac13\sum_{i=1}^3
       \bigl(2+\mathbf1[\text{head}=V,\ \text{tail}_i\ne T]\bigr).
\]
For this cost analysis, represent each equation pair by three weighted
directed ordinary pairs, one from each tail block to the head block.
The baseline two is averaged along with the extra cost. Let $\mathcal P_e$
be the multiset of all these replacement pairs and $\mathcal P_o$ the
multiset of opposite pairs, directed from their tail to their head. Set
$\mathcal P=\mathcal P_o\mathbin{\dot\cup}\mathcal P_e$ and assign
\[
 \omega(p)=1\ (p\in\mathcal P_o),\qquad \omega(p)=1/3\ (p\in\mathcal P_e),\qquad \sum_{p\in\mathcal P}\omega(p)=6m+\frac13(36m)=18m.
\]
These weights describe the cost analysis; the constructed graph
remains unweighted. Write $\operatorname{tail}(p)$ and
$\operatorname{head}(p)$ for the two literal blocks of $p$.
Let $\mathcal B=\{A_x^v:x\in\mathcal X(F),\ v\in\{0,1\}\}$.
In each case below, consider only block-state maps induced by a
cover of the corresponding type.
For a fixed block-state map $z:\mathcal B\to\{T,U,V\}$, its total
literal-block cost is
\[
 L(z)=\sum_{x\in\mathcal X(F)}\sum_{v\in\{0,1\}}
 d_x\left(9\mathbf1[z(A_x^v)=T]
              +\frac92\mathbf1[z(A_x^v)=U]\right).
\]
Put
\[
 I_p(z)=\mathbf1[z(\operatorname{tail}(p))\ne T,\,
                 z(\operatorname{head}(p))=V]
 \qquad(p\in\mathcal P).
\]
Let $C_{\mathrm{np}}(z)$ be the minimum literal-block and non-padding-pair
cost among $B$-first covers inducing $z$. The preceding inequality gives
\begin{equation}\label{eq:np-total-depth-weighted-surrogate}
 C_{\mathrm{np}}(z)
 \overset{\max\geq\text{average}}{\geq}
 L(z)+\sum_{p\in\mathcal P}\omega(p)\bigl(2+I_p(z)\bigr)
 =36m+L(z)+\sum_{p\in\mathcal P}\omega(p)I_p(z).
\end{equation}
For every $x\in\mathcal X(F)$ and $v\in\{0,1\}$,
\begin{equation}\label{eq:np-total-depth-pair-degrees}
 \sum_{p\in\mathcal P:\,\operatorname{head}(p)=A_x^v}\omega(p)
 =
 \sum_{p\in\mathcal P:\,\operatorname{tail}(p)=A_x^v}\omega(p)
 =d_x+2d_x=3d_x.
\end{equation}
Indeed, opposite pairs contribute $d_x$ in each direction. In each
equation containing $x$, the literal is a tail in two satisfying
assignments and all three positions, each with weight $1/3$, and a
head in two assignments at its own position, each with total weight one.

Every directed pair satisfies the pointwise inequality
\[
 1\leq I_p(z)+\mathbf1[z(\operatorname{head}(p))\ne V]
                  +\mathbf1[z(\operatorname{tail}(p))=T].
\]
Indeed, if neither block indicator pays the unit, the head is $V$
and the tail is not $T$, so $I_p(z)=1$. Multiply by $\omega(p)$
and sum, using $\sum_{p\in\mathcal P}\omega(p)=18m$.
The equal weighted in- and out-degrees give the certificate
\begin{equation}\label{eq:np-total-depth-charge-payment}
\begin{gathered}
 18m
 \overset{\eqref{eq:np-total-depth-pair-degrees}}{\leq}
 \sum_{p\in\mathcal P}\omega(p)I_p(z)+\sum_{x\in\mathcal X(F)}\sum_{v\in\{0,1\}}
   3d_x\bigl(2\mathbf1[z(A_x^v)=T]+\mathbf1[z(A_x^v)=U]\bigr)\\
 \overset{6d_x\leq9d_x,\ 3d_x\leq9d_x/2}{\leq}
 \sum_{p\in\mathcal P}\omega(p)I_p(z)+L(z).
\end{gathered}
\end{equation}
For $C$-first covers, use $C_{\mathrm{np}}(z)$ for the corresponding
minimum and reverse all ordinary-pair directions; the same bounds
apply by~\eqref{eq:np-total-depth-pair-degrees}.
For either competing cover, this certificate and the weighted baseline
prove
\begin{equation}\label{eq:np-total-depth-rival-charging}
 C_{\mathrm{np}}(z)
 \overset{\eqref{eq:np-total-depth-weighted-surrogate},\,
          \eqref{eq:np-total-depth-charge-payment}}{\geq}36m+18m=54m.
\end{equation}

Each padding pair costs at least two, since one of its endpoints must
be in $S_0$. Therefore
\begin{equation}\label{eq:np-total-depth-rival-costs}
 \operatorname{cost}_B(G),\operatorname{cost}_C(G)
 \overset{\eqref{eq:np-total-depth-rival-charging}}{\geq}
 3(18m)+2(18m)=90m.
\end{equation}
Applying Lemma~\ref{lem:partial-parity-completion} to the empty
partial assignment gives
\begin{equation}\label{eq:np-total-depth-opt-range}
 \frac m2\leq\OPT(F)\leq m.
\end{equation}
It follows that
\begin{equation}\label{eq:np-total-depth-first-cover-minimum}
 \operatorname{cost}_A(G)
 \overset{\eqref{eq:np-total-depth-A-cost}}{=}93m-6\OPT(F)
  \overset{\eqref{eq:np-total-depth-opt-range}}{\leq}90m
 \overset{\eqref{eq:np-total-depth-rival-costs}}{\leq}
  \min\{\operatorname{cost}_B(G),\operatorname{cost}_C(G)\}.
\end{equation}
Thus the $A$-first cost attains the minimum in the path identity, giving
\[
\begin{aligned}
 W(G)&\overset{\eqref{eq:np-total-depth-path-cost-formula},\,
                    \eqref{eq:np-total-depth-first-cover-minimum}}{=}
       N+3+\operatorname{cost}_A(G)\\
 &\overset{\eqref{eq:np-total-depth-part-sizes},\,\eqref{eq:np-total-depth-A-cost}}{=}
       99m+3+93m-6\OPT(F)=192m-6\OPT(F)+3
  \overset{\eqref{eq:np-total-depth-opt-range}}{\geq}186m+3.
\end{aligned}
\]
\end{proof}

\subsubsection{The approximation gap}

\begin{theorem}[Dense total-depth gap]
\label{thm:np-total-depth-dense-gap}
For every fixed $1\leq r<63/62$, it is NP-hard to approximate total
elimination depth within factor $r$ on connected simple unweighted
graphs. The reduction starts from a bounded-occurrence parity source,
with a fixed error $\epsilon=\epsilon(r)>0$ chosen in the proof.
Its tripartite graph before clique completion satisfies
$\Delta(G_F^\Sigma)\leq14b_\epsilon$.
\end{theorem}

\begin{proof}[Proof sketch]
Proposition~\ref{prop:np-total-depth-parity-identity} gives YES and
NO bounds $(186+6\epsilon)m+3$ and $(189-6\epsilon)m+3$ for
$W(G_F^\Sigma)$. By Lemma~\ref{lem:np-total-depth-added-vertices},
the output total depth lies between $MW(G_F^\Sigma)$ and
$MW(G_F^\Sigma)+N(N+1)/2$, where $N=99m$.
For a fixed $r<189/186=63/62$, choose a small source error,
repeat the source to absorb the additive three, and take $M=C_rm$
with a sufficiently large fixed integer $C_r\geq1$ to absorb the quadratic
remainder. The output still has $O(m)$ vertices and $O(m^2)$ edges.
Its completed core is connected because auxiliary pairs join $B$
to $C$ and literal-block edges join $A$ to both; added vertices
have neighbors in their corresponding parts.
Appendix~\ref{app:dense-np-total} gives all parameter choices and
the full separation calculation.
\end{proof}
